\section{11.09.2026} 

\[
    \Pr(\mathcal{A}) = \frac{N_{\text{благоприятных}} }{N_{\text{общих}}}
\]

\subsection{Операции с множествами}

Будем говорить о счетных множествах

\begin{definition}
    Множество --- совокупность элементов произвольной природы
\end{definition}

\begin{definition}
    Подмножество --- любая совокупность элементов множества, включая пустое множество и само множество
\end{definition}

\begin{definition}
    Пересечение, объединение $\ldots$
\end{definition}

\begin{definition}
    Формула включений исключений выражает количество элементов в объединенном множестве через количество элементов в подмножествах объединенного множества.

    \[
        N(A + B + C) = N(A) + N(B) + N(C) - N(AB) - N(BC) - N(AC) + N(ABC)
    \]

    Общий вид $\ldots$
\end{definition}

\begin{exercise}[Задача со вступительных в Яндекс]
На концерте каждую песню исполняли двое артистов, и никакая пара не выступала вместе более одного раза. Всего было исполнено 30 песен, и каждый артист выступил по 5 раз. Сколько всего выступало артистов?
\end{exercise}
\begin{solution}
$\ldots$

\[
    30 = 5n - 30 \xrightarrow[]{} n = 12
\]
\end{solution}

\subsection{Классическая вероятность}

\begin{definition}
\end{definition}

\begin{exercise}
    Какова вероятность того, что при шести бросаниях правильной шестигранной игральной кости хотя бы один раз выпадет шесть очков?
\end{exercise}
\begin{solution}
    Интуитивно кажется, что вероятность близка к $100\%$ --- давайте ее вычислим.

    Сначала нужно начертить дерево решений - визуализировать

    \[
        \Pr_{}(\mathcal{A} ) = \frac{}{6^6} = 1 - \frac{\overline{N_{\text{благоприятных}}} }{N_{\text{общих}}} = \Pr_{}(\overline{\mathcal{A}} ) \text{ --- ни разу не выпала 6} = 1 - \left( \frac{5}{6} \right)^6 \approx 67\% 
    \]
\end{solution}

\subsection{Основы комбинаторики}

\begin{definition}
    Перестановки, сочетания, размещения
\end{definition}

Использование формул должно быть аргументировано (например схематически)

\subsection{Задачи}

\subsubsection{Задача 1}

\begin{exercise}
Ребенок играет с 10 буквами разрезной азбуки, А, А, А, Е, И, К, М, М, Т, Т. Какова вероятность, что при случайном расположении букв в ряд он получит слово «математика»?
\end{exercise}

\begin{solution}
\[
    \text{М А Т Е М А Т И К А}
\]

    Смотрим дерево 

    Выпишем формулу классической вероятности:

\[
    \Pr(\mathcal{A}) = \frac{N_{\text{благоприятных}} }{N_{\text{общих}}} = \frac{3! \cdot 2! \cdot 2! \cdot 1! \cdot 1! \cdot 1!}{10!}
\]

Важно: мы должны показывать в следствиях что исходы равновероятны (проверять). На коллоквиуме можно показать на примере дерева решений, только после этого можно считать вероятность.
\end{solution}

\subsubsection{Задача 2}

\begin{exercise}
В конце 2004 года компания МТС предлагала своим новым абонентам следующую лотерею. Выдавался билет с 25 закрашенными клетками ($5\times5$), в каждой клетке "скрывалась" буква. В каждой строке все буквы были различны. Необходимо было стереть по одной клетке в каждой строке и получить слово "ДЖИНС" (в правильном порядке, т.е. "Д" из первой строки, "Ж" из второй и т.д.). Игра была честная, т.е. в каждой строке была необходимая буква. Найти вероятность выигрыша.
\end{exercise}

\begin{solution}
    Смотрим дерево 

    Выпишем формулу классической вероятности:

\[
    \Pr(\mathcal{A}) = \frac{N_{\text{благоприятных}} }{N_{\text{общих}}} = \frac{1}{5 \cdot 5 \cdot 5 \cdot 5 \cdot 5}
\]
\end{solution}
